% Part: second-order-logic
% Chapter: syntax-and-semantics
% Section: language-of-sol

\documentclass[../../../include/open-logic-section]{subfiles}

\begin{document}

\olfileid{sol}{syn}{lan}

\olsection{في لغة منطق الرتبة الثانية}

تُنشأ عبارات منطق الرتبة الثانية، كما تُنشأ عبارات الرتبة الأولى، من
مفردات أصلية: \emph{!!{variable}s} و\emph{!!{constant}s}
و\emph{!!{predicate}s}، وربما \emph{!!{function}s}. فإذا ضممنا إليها
الروابط المنطقية والمكمّمات وعلامات الفصل، من أقواس وفواصل، أنشأنا
\emph{الحدود} و\emph{!!{formula}s}. وإنما يمتاز منطق الرتبة الثانية بأنه
يزيد على متغيرات الأفراد متغيرات للعلاقات والدوال، ويجيز التكميم عليها.
فرموزه المنطقية هي رموز منطق الرتبة الأولى مع الصنفين الآتيين:

\begin{enumerate}
\item مجموعة !!{denumerable} من !!{variable}s العلاقات في الرتبة الثانية،
  لكل رتبة~$n$: $\Obj V_0^n$، $\Obj V_1^n$، $\Obj V_2^n$، \dots
\item ومجموعة !!{denumerable} من !!{variable}s الدوال في الرتبة الثانية،
  لكل رتبة~$n$: $\Obj u_0^n$، $\Obj u_1^n$، $\Obj u_2^n$، \dots
\end{enumerate}

ويُدخل في منطق الرتبة الأولى عادةً رمز علاقة !!{identity}~$\eq$. ثم إن
الرموز غير المنطقية في لغة~$\Lang{L}$ تكاد تكون لازمة للتعبير عن معنى
ذي بال. نعم، توجد !!{sentence}s لا تستعمل رمزًا غير منطقي، ولكن لا يكاد
يقال، باستعمال~$\eq$ وحدها، شيء ذو شأن. وأما منطق الرتبة الثانية، ففيه
من متغيرات العلاقات والدوال عدد غير محدود؛ ولذلك نستطيع، من غير مخزون
خاص من الرموز غير المنطقية، أن نقول كل ما تقوله لغة من الرتبة الأولى.

\end{document}
